\documentclass[a4paper,11pt]{ltjsarticle}
\usepackage{luatexja-fontspec}\usepackage[haranoaji]{luatexja-preset}
\usepackage{amsmath,amssymb,amsthm,bm,mathtools}
\usepackage{booktabs,array}\usepackage{geometry}\usepackage{xcolor}
\usepackage{pgfplots}\pgfplotsset{compat=1.18}
\geometry{margin=24mm}
\newcommand{\tr}{\operatorname{tr}}\newcommand{\rank}{\operatorname{rank}}
\newcommand{\OO}{\mathcal{O}}
\title{K-recursion Gaussian BN: 応用実験サマリー\\\large 隠れノード推定・センサー配置・介入の情報幾何・構造学習}
\author{自動生成 (experiments/make\_summary.py より)}
\date{2026年6月10日}
\begin{document}\maketitle
\begin{abstract}
\texttt{gaussian\_bn} ライブラリを用いた3つの応用実験の\textbf{実行結果}である。全数値は各 \texttt{run} スクリプトが出力した \texttt{results/exp\{1,2,3\}.json} から機械転記しており、合成データは含まない。乱数シードは \texttt{20260610}。実験はノートの Experiment 2--4 と Theme D に対応する。
\end{abstract}
\section{実験1: 隠れノード推定 --- EM と勾配学習の比較}
diamond DAG（$0\to1,0\to2,1\to3,2\to3$）の中間 $\{1,2\}$ を hidden とし、観測 $\OO=\{0,3\}$（$N=200000$）。観測周辺尤度 $\mathcal L=\log\det K_{\OO\OO}+\tr(K_{\OO\OO}^{-1}S)$ を EM・Adam・LBFGS で最小化した。
\begin{center}\begin{tikzpicture}
\begin{semilogyaxis}[width=0.78\linewidth,height=6cm,xlabel={反復},ylabel={NLL $-$ 最適値},grid=both,legend pos=north east]
\addplot[blue,thick] coordinates {(0,0.272813) (1,0.02189) (2,0.00406684) (3,0.000768936) (4,0.00014586) (5,2.76856e-05) (6,5.25559e-06) (7,9.97702e-07) (8,1.89401e-07) (9,3.59554e-08) (10,6.82568e-09) (11,1.29577e-09) (12,2.45986e-10) (13,4.66973e-11) (14,8.86491e-12) (15,1.6831e-12) (16,3.193e-13) (17,6.08402e-14) (18,1.15463e-14) (19,2.22045e-15) (20,4.44089e-16) (21,1e-16) (22,1e-16) (23,1e-16) (24,1e-16) (25,4.44089e-16) (26,1e-16) (27,4.44089e-16) (28,1e-16) (29,1e-16) (30,1e-16) (31,1e-16) (32,1e-16) (33,1e-16) (34,1e-16) (35,4.44089e-16) (36,1e-16) (37,1e-16) (38,1e-16) (39,1e-16) (40,1e-16) (41,1e-16) (42,4.44089e-16) (43,1e-16) (44,1e-16) (45,4.44089e-16) (46,4.44089e-16) (47,4.44089e-16) (48,4.44089e-16) (49,4.44089e-16) (50,1e-16) (51,4.44089e-16) (52,4.44089e-16) (53,1e-16) (54,1e-16) (55,1e-16) (56,1e-16) (57,1e-16) (58,1e-16) (59,1e-16)};
\addplot[red,thick] coordinates {(0,0.883917) (1,0.684644) (2,0.507225) (3,0.355956) (4,0.235046) (5,0.147167) (6,0.0915406) (7,0.0641798) (8,0.0589069) (9,0.0668342) (10,0.0778745) (11,0.084151) (12,0.0826795) (13,0.075794) (14,0.0682859) (15,0.0629242) (16,0.0585673) (17,0.052526) (18,0.0435548) (19,0.0325358) (20,0.0215904) (21,0.0129147) (22,0.00794054) (23,0.00698713) (24,0.00934922) (25,0.0136591) (26,0.0183367) (27,0.0219847) (28,0.0236522) (29,0.0229559) (30,0.0200763) (31,0.0156574) (32,0.0106328) (33,0.00600451) (34,0.00260573) (35,0.000898675) (36,0.000866405) (37,0.00204864) (38,0.00372613) (39,0.0051881) (40,0.00596745) (41,0.00594283) (42,0.00528334) (43,0.00429542) (44,0.00326687) (45,0.00237828) (46,0.00169661) (47,0.00121925) (48,0.000923343) (49,0.000791265) (50,0.000808769) (51,0.000949654) (52,0.00116372) (53,0.00137831) (54,0.00151431) (55,0.00151043) (56,0.00134525) (57,0.00104754) (58,0.000689075) (59,0.000361151) (60,0.000142765) (61,7.28476e-05) (62,0.000137997) (63,0.000281151) (64,0.000427255) (65,0.000514017) (66,0.000513743) (67,0.000437461) (68,0.000322053) (69,0.000209276) (70,0.000127949) (71,8.66609e-05) (72,7.74936e-05) (73,8.5681e-05) (74,9.86098e-05) (75,0.000109905) (76,0.000118276) (77,0.000123785) (78,0.000124932) (79,0.000118494) (80,0.00010187) (81,7.59332e-05) (82,4.61766e-05) (83,2.10337e-05) (84,8.09883e-06) (85,1.03065e-05) (86,2.42767e-05) (87,4.18463e-05) (88,5.39992e-05) (89,5.50702e-05) (90,4.50177e-05) (91,2.87482e-05) (92,1.31443e-05) (93,3.59248e-06) (94,1.82287e-06) (95,5.90261e-06) (96,1.19648e-05) (97,1.64463e-05) (98,1.7608e-05) (99,1.57564e-05) (100,1.24042e-05) (101,9.10346e-06) (102,6.6731e-06) (103,5.1243e-06) (104,4.09544e-06) (105,3.34777e-06) (106,2.95364e-06) (107,3.11451e-06) (108,3.83465e-06) (109,4.75391e-06) (110,5.28847e-06) (111,4.97967e-06) (112,3.7973e-06) (113,2.1822e-06) (114,7.99683e-07) (115,1.676e-07) (116,3.91538e-07) (117,1.15361e-06) (118,1.93301e-06) (119,2.30275e-06) (120,2.12576e-06) (121,1.56044e-06) (122,9.11134e-07) (123,4.39708e-07) (124,2.5046e-07) (125,2.91254e-07) (126,4.35272e-07) (127,5.70866e-07) (128,6.45028e-07) (129,6.53482e-07) (130,6.0809e-07) (131,5.16744e-07) (132,3.87672e-07) (133,2.43841e-07) (134,1.25454e-07) (135,7.22926e-08) (136,9.79845e-08) (137,1.77225e-07) (138,2.57784e-07) (139,2.90365e-07) (140,2.56262e-07) (141,1.749e-07) (142,8.80517e-08) (143,3.31574e-08) (144,2.36724e-08) (145,4.71594e-08) (146,7.86645e-08) (147,9.77157e-08) (148,9.764e-08) (149,8.3729e-08) (150,6.50819e-08) (151,4.77663e-08) (152,3.35146e-08) (153,2.25594e-08) (154,1.62586e-08) (155,1.6621e-08) (156,2.35703e-08) (157,3.31608e-08) (158,3.90533e-08) (159,3.65053e-08) (160,2.5767e-08) (161,1.21187e-08) (162,2.41559e-09) (163,6.79788e-10) (164,5.86964e-09) (165,1.3163e-08) (166,1.74835e-08) (167,1.66051e-08) (168,1.18603e-08) (169,6.45832e-09) (170,3.04419e-09) (171,2.29292e-09) (172,3.18947e-09) (173,4.30338e-09) (174,4.85419e-09) (175,4.87522e-09) (176,4.68371e-09) (177,4.35755e-09) (178,3.70407e-09) (179,2.62488e-09) (180,1.4101e-09) (181,6.21539e-10) (182,6.57441e-10) (183,1.39922e-09) (184,2.26173e-09) (185,2.60944e-09) (186,2.19741e-09) (187,1.29913e-09) (188,4.57383e-10) (189,8.71929e-11) (190,2.35334e-10) (191,6.34559e-10) (192,9.5086e-10) (193,1.00689e-09) (194,8.35912e-10) (195,5.82169e-10) (196,3.69354e-10) (197,2.41387e-10) (198,1.86818e-10) (199,1.86977e-10) (200,2.32113e-10) (201,3.03538e-10) (202,3.59684e-10) (203,3.53978e-10) (204,2.71833e-10) (205,1.49279e-10) (206,5.11791e-11) (207,2.48921e-11) (208,6.78235e-11) (209,1.34214e-10) (210,1.72355e-10) (211,1.60099e-10) (212,1.12594e-10) (213,6.23843e-11) (214,3.31473e-11) (215,2.79745e-11) (216,3.60814e-11) (217,4.61702e-11) (218,5.33267e-11) (219,5.66787e-11) (220,5.46945e-11) (221,4.50635e-11) (222,2.87685e-11) (223,1.23919e-11) (224,4.54303e-12) (225,8.98659e-12) (226,2.09766e-11) (227,3.06186e-11) (228,3.0477e-11) (229,2.09037e-11) (230,8.84315e-12) (231,1.7657e-12) (232,2.22178e-12) (233,7.14051e-12) (234,1.14593e-11) (235,1.22506e-11) (236,1.00249e-11) (237,7.02594e-12) (238,4.86677e-12) (239,3.67617e-12) (240,2.93054e-12) (241,2.502e-12) (242,2.72316e-12) (243,3.65219e-12) (244,4.60476e-12) (245,4.61275e-12) (246,3.34177e-12) (247,1.49436e-12) (248,2.53131e-13) (249,2.97096e-13) (250,1.29363e-12) (251,2.27596e-12) (252,2.47491e-12) (253,1.85363e-12) (254,9.57456e-13) (255,3.6815e-13) (256,2.76668e-13) (257,4.96492e-13) (258,7.34968e-13) (259,8.40217e-13) (260,8.15792e-13) (261,7.0699e-13) (262,5.37348e-13) (263,3.32623e-13) (264,1.64757e-13) (265,1.19016e-13) (266,2.13607e-13) (267,3.62377e-13) (268,4.35652e-13) (269,3.7037e-13) (270,2.14939e-13) (271,7.72715e-14) (272,3.5083e-14) (273,7.9492e-14) (274,1.46105e-13) (275,1.78968e-13) (276,1.63869e-13) (277,1.25233e-13) (278,8.57092e-14) (279,5.55112e-14) (280,3.73035e-14) (281,3.33067e-14) (282,4.66294e-14) (283,6.88338e-14) (284,8.12683e-14) (285,7.01661e-14) (286,4.04121e-14) (287,1.11022e-14) (288,2.66454e-15) (289,1.59872e-14) (290,3.4639e-14) (291,4.21885e-14) (292,3.4639e-14) (293,1.9984e-14) (294,9.32587e-15) (295,7.10543e-15) (296,1.02141e-14) (297,1.33227e-14) (298,1.42109e-14) (299,1.37668e-14) (300,1.33227e-14) (301,1.11022e-14) (302,7.54952e-15) (303,3.9968e-15) (304,2.66454e-15) (305,4.88498e-15) (306,7.54952e-15) (307,9.32587e-15) (308,7.54952e-15) (309,3.9968e-15) (310,1.33227e-15) (311,1.77636e-15) (312,3.55271e-15) (313,3.9968e-15) (314,4.44089e-15) (315,3.55271e-15) (316,3.10862e-15) (317,1.77636e-15) (318,1.77636e-15) (319,1.77636e-15) (320,2.22045e-15) (321,2.22045e-15) (322,2.66454e-15) (323,2.66454e-15) (324,1.77636e-15) (325,1.33227e-15) (326,1.33227e-15) (327,1.77636e-15) (328,1.77636e-15) (329,1.77636e-15) (330,1.77636e-15) (331,1.77636e-15) (332,8.88178e-16) (333,1.77636e-15) (334,1.33227e-15) (335,1.33227e-15) (336,1.33227e-15) (337,1.77636e-15) (338,1.33227e-15) (339,8.88178e-16) (340,1.77636e-15) (341,8.88178e-16) (342,1.33227e-15) (343,1.33227e-15) (344,1.77636e-15) (345,1.33227e-15) (346,1.33227e-15) (347,1.33227e-15) (348,1.77636e-15) (349,1.77636e-15) (350,1.33227e-15) (351,1.33227e-15) (352,1.33227e-15) (353,1.77636e-15) (354,1.33227e-15) (355,1.33227e-15) (356,1.33227e-15) (357,8.88178e-16) (358,8.88178e-16) (359,1.33227e-15) (360,8.88178e-16) (361,1.77636e-15) (362,1.33227e-15) (363,8.88178e-16) (364,8.88178e-16) (365,1.33227e-15) (366,1.33227e-15) (367,1.33227e-15) (368,1.33227e-15) (369,8.88178e-16) (370,8.88178e-16) (371,1.33227e-15) (372,1.33227e-15) (373,1.33227e-15) (374,1.33227e-15) (375,8.88178e-16) (376,1.33227e-15) (377,8.88178e-16) (378,8.88178e-16) (379,1.33227e-15) (380,1.33227e-15) (381,8.88178e-16) (382,1.77636e-15) (383,8.88178e-16) (384,8.88178e-16) (385,8.88178e-16) (386,1.33227e-15) (387,8.88178e-16) (388,1.33227e-15) (389,1.33227e-15) (390,8.88178e-16) (391,1.33227e-15) (392,8.88178e-16) (393,8.88178e-16) (394,1.33227e-15) (395,8.88178e-16) (396,1.33227e-15) (397,8.88178e-16) (398,1.33227e-15) (399,1.77636e-15)};
\legend{EM,Adam}\end{semilogyaxis}\end{tikzpicture}\end{center}
\begin{center}\begin{tabular}{lcccc}\toprule
手法 & 最終 NLL & $\|K_{\OO\OO}(\hat\eta)-S\|/\|S\|$ & 反復 & 時間[s] \\ \midrule
EM & 2.203 & $2.327\times 10^{-15}$ & 60 & 0.220 \\
Adam & 2.203 & $4.880\times 10^{-10}$ & 400 & 1.337 \\
LBFGS & 2.203 & $1.696\times 10^{-9}$ & 9 & 0.023 \\
\bottomrule\end{tabular}\end{center}
NLL の理論最適値は $\log\det S+\dim=2.203$ で、3手法すべてが到達した（観測共分散を標本共分散へ復元）。
\paragraph{gauge 診断.}解での4 edge 母数の Fisher 固有値は $\{-2.541\times 10^{-16},\ 1.863\times 10^{-16},\ 0.280,\ 0.965\}$、$\rank G^{(\OO)}=2<q=4$。零空間は2次元で、2つの hidden ノードのスケール gauge に対応する（観測共分散は復元できても edge 母数自体は識別不能）。
\section{実験2: センサー配置と edge 識別可能性}
6ノード DAG の6本の edge 母数に対し、観測ノード集合 $\OO$ を予算 $m$ 以内で選ぶ。pullback Fisher 計量 $G^{(\OO)}$ の D最適 $\log\det(G^{(\OO)}+\epsilon I)$ を greedy 最大化し、各予算での識別 rank と最小固有値も記録した。
\begin{center}\begin{tikzpicture}
\begin{axis}[width=0.78\linewidth,height=5.6cm,xlabel={観測ノード数(予算 $m$)},ylabel={$\log\det G^{(\OO)}$},axis y line*=left,grid=both]
\addplot[blue,mark=*,thick] coordinates {(1,-102.84) (2,-60.5488) (3,-1.99724) (4,1.93769) (5,4.96084) (6,6.22977)};\label{plt:ld}
\end{axis}
\begin{axis}[width=0.78\linewidth,height=5.6cm,axis y line*=right,ylabel={識別 rank},ymin=0,ymax=6.5,hide x axis]
\addlegendimage{blue,mark=*}\addlegendentry{$\log\det G$}
\addplot[red,mark=square,thick,dashed] coordinates {(1,1) (2,3) (3,6) (4,6) (5,6) (6,6)};\addlegendentry{rank}
\end{axis}\end{tikzpicture}\end{center}
\begin{center}\begin{tabular}{ccccc}\toprule
$m$ & 選択ノード & $\log\det G$ & $\lambda_{\min}$ & rank \\ \midrule
1 & [5] & $-102.840$ & $1.000\times 10^{-9}$ & 1/6 \\
2 & [5, 3] & $-60.549$ & $1.000\times 10^{-9}$ & 3/6 \\
3 & [5, 3, 2] & $-1.997$ & $0.051$ & 6/6 \\
4 & [5, 3, 2, 1] & $1.938$ & $0.324$ & 6/6 \\
5 & [5, 3, 2, 1, 4] & $4.961$ & $0.855$ & 6/6 \\
6 & [5, 3, 2, 1, 4, 0] & $6.230$ & $1.650$ & 6/6 \\
\bottomrule\end{tabular}\end{center}
わずか 3 ノードの観測で全6 edge が識別可能（rank $=6$）になり、以降はセンサー追加で条件数（$\log\det G,\ \lambda_{\min}$）が改善する。
\paragraph{greedy vs exhaustive (予算3).}D最適では greedy [5, 3, 2] と exhaustive [2, 3, 5] が一致（スコア -1.997）。一方 E最適（$\max\lambda_{\min}$）は greedy [0, 1, 2] ($1.000\times 10^{-9}$) が exhaustive [1, 4, 5] ($0.066$) に劣る。これは E最適性が劣モジュラでないための既知の挙動で、rank 不足の領域では greedy が停滞する（exhaustive か rank 考慮基準が必要）。
\section{実験3: 介入の情報幾何}
交絡三角形 $0\to1,\ 0\to2,\ 1\to2$（ノード0は隠れ交絡、$1\to2$ が直接効果 $c$）。観測上の連関 $I(V_1;V_2)$ は交絡経路と直接経路の混合だが、hard 介入 $\mathrm{do}(V_1)$ はノード1を親から切り離すため $I^{\mathrm{do}}(V_1;V_2)$ は直接因果効果のみを表す。差 $\Delta I=I(V_1;V_2)-I^{\mathrm{do}}(V_1;V_2)$ が交絡（非因果）成分である。
\begin{center}\begin{tikzpicture}
\begin{axis}[width=0.78\linewidth,height=5.8cm,xlabel={直接効果 $c$},ylabel={相互情報量 [nats]},grid=both,legend pos=north west]
\addplot[blue,mark=*,thick] coordinates {(0.0,0.377715) (0.15,0.469627) (0.3,0.558674) (0.45,0.643862) (0.6,0.724804) (0.75,0.801456) (0.9,0.873959) (1.05,0.942542) (1.2,1.00747) (1.35,1.06902) (1.5,1.12746)};
\addplot[red,mark=square,thick] coordinates {(0.0,0) (0.15,0.00231422) (0.3,0.00919331) (0.45,0.0204522) (0.6,0.0358006) (0.75,0.0548662) (0.9,0.0772226) (1.05,0.102416) (1.2,0.129991) (1.35,0.159508) (1.5,0.190558)};
\legend{$I(V_1;V_2)$ (観測),$I^{\mathrm{do}}(V_1;V_2)$ (介入)}
\end{axis}\end{tikzpicture}\end{center}
\begin{center}\begin{tabular}{lccc}\toprule
設定 & $I(V_1;V_2)$ & $I^{\mathrm{do}}(V_1;V_2)$ & 交絡成分 $\Delta I$ \\ \midrule
$c=0$（交絡のみ） & 0.378 & $0$ & 0.378 \\
$c=0.9$（交絡+直接） & 0.874 & 0.077 & 0.797 \\
\bottomrule\end{tabular}\end{center}
$c=0$ では観測連関があるのに $I^{\mathrm{do}}=0$（連関は完全に交絡由来；介入で$\mathrm{Cov}(V_1,V_2)$ も $0$）。$c$ が増えると $I^{\mathrm{do}}$ が直接因果効果として立ち上がり、観測連関との差が交絡成分を与える。do 演算が因果と交絡を情報量で分離できる。
\section{実験5: 構造学習（group sparsity による枝刈り）}
全結合 supergraph（$i<j$ の全10本）から出発し、5ノード（各2次元）の BN を K-recursion 周辺尤度＋\textbf{group sparsity} で学習：
\begin{equation}\min_\eta\ \mathcal L(\eta)+\lambda\sum_{(i,j)}\lVert A_{ji}\rVert_F\end{equation}
を近接勾配（平滑勾配ステップ＋edge ブロックの group soft-threshold）で解く。$\lambda$ を上げると不要 edge の $2\times2$ ブロックが丸ごと $0$ に潰れる。各 $\lambda$ の選択構造を無罰則 MLE（\texttt{fit\_local\_regression}）で再フィットし BIC で構造選択した。
\begin{center}\begin{tikzpicture}
\begin{axis}[width=0.78\linewidth,height=5.8cm,xlabel={正則化 $\lambda$},ylabel={edge ブロック $\lVert A_{ji}\rVert_F$},grid=both,legend pos=north east]
\addplot[blue,thick,mark=*] coordinates {(0.0,1.097) (1.0,0.794381) (2.0,0) (4.0,0) (7.0,0)};
\addplot[red,dashed,mark=square,mark size=1pt] coordinates {(0.0,0.128994) (1.0,0) (2.0,0) (4.0,0) (7.0,0)};
\addplot[blue,thick,mark=*] coordinates {(0.0,0.922643) (1.0,0.617311) (2.0,0.29293) (4.0,0) (7.0,0)};
\addplot[blue,thick,mark=*] coordinates {(0.0,0.817475) (1.0,0.48374) (2.0,0) (4.0,0) (7.0,0)};
\addplot[red,dashed,mark=square,mark size=1pt] coordinates {(0.0,0.109019) (1.0,0) (2.0,0) (4.0,0) (7.0,0)};
\addplot[red,dashed,mark=square,mark size=1pt] coordinates {(0.0,0.116895) (1.0,0) (2.0,0) (4.0,0) (7.0,0)};
\addplot[red,dashed,mark=square,mark size=1pt] coordinates {(0.0,0.0839098) (1.0,0) (2.0,0) (4.0,0) (7.0,0)};
\addplot[red,dashed,mark=square,mark size=1pt] coordinates {(0.0,0.117565) (1.0,0) (2.0,0) (4.0,0) (7.0,0)};
\addplot[blue,thick,mark=*] coordinates {(0.0,0.789728) (1.0,0.410796) (2.0,0) (4.0,0) (7.0,0)};
\addplot[blue,thick,mark=*] coordinates {(0.0,1.07951) (1.0,0.741037) (2.0,0) (4.0,0) (7.0,0)};
\addlegendimage{blue,thick,mark=*}\addlegendentry{真の edge (5)}
\addlegendimage{red,dashed,mark=square}\addlegendentry{偽の edge (5)}
\end{axis}\end{tikzpicture}\end{center}
\begin{center}\begin{tabular}{ccccc}\toprule
$\lambda$ & 再フィット NLL & BIC & edge数 & 選択 \\ \midrule
0 & 3.071 & 1557.856 & 10 &  \\
1.000 & 3.140 & 1465.588 & 5 & $\leftarrow$ \\
2.000 & 7.397 & 3072.541 & 1 &  \\
4.000 & 8.539 & 3505.622 & 0 &  \\
7.000 & 8.539 & 3505.622 & 0 &  \\
\bottomrule\end{tabular}\end{center}
$\lambda=0$（罰則なし MLE）は全10本を残して過適合するが、$\lambda=1.000$ で偽 edge の $2\times2$ ブロックが正確に $0$ となり、BIC が真の構造を選択。\textbf{完全復元}（TP=5, FP=0, FN=0, precision $=1.000$, recall $=1.000$）。全観測のため再フィットは局所回帰の閉形式 MLE で行えるが、同じ枠組みは観測集合を変えるだけで hidden node にも拡張できる。
\section{まとめ}
\begin{itemize}
\item \textbf{実験1}: 隠れノードありでも EM/Adam/LBFGS が同一最適 NLL に到達し観測共分散を復元、Fisher rank が latent gauge を正しく検出。
\item \textbf{実験2}: Fisher 情報に基づく D最適 greedy センサー配置で、少数観測でも全 edge を識別可能にできる。E最適は劣モジュラでないため greedy の限界も定量的に確認。
\item \textbf{実験3}: do 演算により観測連関を因果（直接）成分と交絡成分へ情報量で分解。
\item \textbf{実験5}: supergraph から group sparsity 近接勾配で枝刈りし、再フィット＋BIC で真の DAG 構造を完全復元（構造学習への入口）。
\item いずれも \texttt{gaussian\_bn} の単一 API（K-recursion・推論・訓練・Fisher・介入・設計）上で完結し、全数値は実行結果ファイル由来である。
\end{itemize}
\end{document}
